\documentclass[11pt]{article}

\usepackage[T1]{fontenc}
\usepackage{lmodern}
\usepackage{xcolor}
\pagecolor{white}
\usepackage{eso-pic}
\AddToShipoutPictureBG{\AtPageLowerLeft{\color{white}\rule{\paperwidth}{\paperheight}}}

\usepackage[margin=1in]{geometry}
\usepackage{amsmath,amssymb,amsthm}
\usepackage{booktabs}
\usepackage{microtype}

\newtheorem{principle}{Principle}
\newtheorem{remark}{Remark}

\newcommand{\bthree}{b_3}
\newcommand{\flux}{F}
\newcommand{\hh}{H}

\title{Mitchell's Principle of Constraint-First Unification:\\ relations, never values}
\author{Rob Mitchell}
\date{}

\begin{document}
\maketitle

\begin{center}\itshape
``Don't ask a complex system for its values --- ask for its constraints, its count, and its
deficit. Everything else is covariance, not knowledge.''
\end{center}

\begin{abstract}
The unification of string theory is not a formula for the Standard Model's parameters; it is a
complete characterization. We state, and assemble the evidence for, a single principle: a theory
of everything predicts \emph{relations}, never \emph{values}. The exact content of the flux sector
is exhausted by one constraint --- the D3 tadpole, the unique duality invariant --- together with
the topological labels and the counting (nineteen, the Standard Model orbit-space dimension).
Everything else transforms by a covariance table, Mitchell's Laws of Duality Covariance, measured
on supersymmetric vacua at $h^{1,2}=2$ and $h^{1,2}=5$. And there is no selection: not exact (a
no-go theorem), not statistical (the flux measure is flat), not cosmological (the cosmological
constant is locally flat, and the uplift that would supply it is contested). The principle
reduces to one equation: the invariant ring of the flux is $\mathbb{C}[N]$ with $N=\flux^\top
J\hh$. We state the four falsifiers that would kill it.
\end{abstract}

\section{The principle, in one equation}

The whole of the programme follows from a single line of invariant theory. Let $(\flux,\hh)$ be
the type IIB three-form flux pair, valued in $H^3(X,\mathbb{Z})\times H^3(X,\mathbb{Z})$, and let
$J$ be the symplectic intersection form. The duality group $\mathrm{Sp}(\bthree,\mathbb{Z})\times
\mathrm{SL}(2,\mathbb{Z})$ acts on the pair, and the First Fundamental Theorem of invariant theory
(Weyl 1939), specialized to two vectors, gives

\begin{equation}
\mathbb{C}[\flux,\hh]^{\,\mathrm{Sp}(\bthree,\mathbb{Z})\times \mathrm{SL}(2,\mathbb{Z})}
\;=\; \mathbb{C}[N],
\qquad N = \flux^\top J\,\hh .
\label{eq:ring}
\end{equation}

The ring of \emph{exact, duality-invariant} quantities in the flux is generated by a single
quadratic invariant --- the D3 tadpole. Everything this note reports is a consequence of
Eq.~(\ref{eq:ring}), stated in three tiers and never confused: \textbf{proven} (symbolic or
machine-checked), \textbf{measured} (numerical, with a control that could have failed), and
\textbf{conjectured} (open).

\begin{principle}[Constraint-First Unification]
\label{pr:main}
A theory of everything predicts relations, never values. The exact content of a sector is
exhausted by three types --- constraints, labels, and counting. Everything else transforms by a
covariance table. And there is no selection: not exact, not statistical, not cosmological.
\end{principle}

\section{The exact structure is constraint + label + counting}

The theorem tower assembles Eq.~(\ref{eq:ring}) and its consequences; the full proofs and the
symbolic verification are in the companion paper \emph{Constraint-First Theorem}
(cited as paper~1). The chain is:

\begin{itemize}
\item \textbf{Mitchell's Uniqueness.} $N=\flux^\top J\hh$ is the \emph{unique} bilinear symplectic
invariant of the flux, for any compactification: $\flux^\top\flux$, $\hh^\top\hh$, $\flux^\top\hh$
are gauge-dependent, and $\flux^\top J\flux=\hh^\top J\hh=0$.
\item \textbf{Constraint-First.} Every duality-invariant exact polynomial in $(\flux,\hh)$ is a
power of $N$ --- the invariant ring is generated by the tadpole alone.
\item \textbf{Completeness.} Over $\mathbb{R}$ the duality group acts transitively on each nonzero
tadpole level set; over $\mathbb{Z}$ the orbits are classified by the content $d=\gcd(\flux,\hh)$
with $d^2\mid N$, finitely many per tadpole.
\item \textbf{Full-duality.} The statement holds for the full group
$\mathrm{Sp}(\bthree,\mathbb{Z})\times\mathrm{SL}(2,\mathbb{Z})$, S- and T-duality included.
\item \textbf{Complete Exact Structure.} The exact content is \emph{exhausted} by constraints (the
tadpole $N\le L$, flux quantization, RR tadpoles, anomaly cancellation) and labels (the topology
$h^{1,1},h^{2,1},c_2,\kappa$). The superpotential $W_0$, the string coupling $g_s$, the masses,
the couplings, and the cosmological constant are provably \emph{not} exact observables.
\end{itemize}

The immediate consequence is a no-go: \emph{no exact, polynomial, duality-invariant selector can
single out a vacuum}. This rules out the equation; it does not rule out statistical,
moduli-dependent, or non-polynomial selection --- a distinction the later sections take up.

\section{The covariance table}

The theorem says what is invariant; the laws say what is covariant. Measured on $100$ primitive
supersymmetric flux vacua at $h^{1,2}=2$ and $12$ in-cone supersymmetric vacua at $h^{1,2}=5$,
every measurement gated by a control, the transformation of the surviving observables under
$\mathrm{SL}(2,\mathbb{Z})$ is:

\begin{center}
\begin{tabular}{lll}
\toprule
law & content & $h^{1,2}=5$ \\
\midrule
I (scaling) & $W(d\flux,d\hh)=d\,W$, $N\to d^2N$ & symbolic \\
II (magnitude) & $|W_S|/|W|=1$ & $1.000000$ (std $3.7\times10^{-15}$) \\
IIb (phase) & $\arg(W_S/W)=-\arg\tau$ & residual $7.4\times10^{-16}$ \\
III ($g_s$) & $g_s\to|\tau|^2 g_s$ & $1.000000$ (std $1.8\times10^{-16}$) \\
IV (moduli) & complex-structure moduli inert at degree $0$ & --- \\
V (counting) & $19=$ SM orbit-space dimension & --- \\
\bottomrule
\end{tabular}
\end{center}

These are the textbook type IIB S-duality laws, exactly. The laws are $h^{1,2}$-independent: the
covariance table does not scatter as moduli are added. In parallel, the supersymmetric-vacuum
fraction collapses with the number of moduli kept --- the in-cone fraction falls $80\%\to
25\%\to 0\%$ from $h^{1,2}=2$ to $5$ to $10$, while the F-flat points still exist. The physical
vacua move \emph{out} of the K\"ahler cone; the exact structure does not move. The collapse is a
landscape fact, not a proposal artefact.

\section{The Standard-Model frontier}

One sector up, the same decomposition holds. The Standard Model's exact content is constraints
(anomaly cancellation, proven $5/5$; hypercharge uniqueness, proven $4/4$), labels (the gauge
group $\mathrm{SU}(3)\times\mathrm{SU}(2)\times\mathrm{U}(1)$, $N_{\rm gen}=|\chi|/2$ a topological
index), and parameters. The flavor group $\mathrm{U}(3)^5$ leaves an orbit space of dimension
nineteen (Molien--Weyl), matching the literature count exactly: six quark masses, three mixing
angles, one phase, three charged-lepton masses, three gauge couplings, two Higgs parameters, and
$\theta_{\rm QCD}$. Ultraviolet completions do impose exact relations --- minimal SU(5)
collapses five (gauge $2$ plus the three $Y_e=Y_d^{\top}$ mass relations $m_e=m_d$, $m_\mu=m_s$,
$m_\tau=m_b$), and minimal SO(10), $E_6$, Pati--Salam, and trinification collapse twelve
($Y_u=Y_d=Y_e$, $V_{\rm CKM}=1$) --- but those relations are \emph{numerically wrong} against
data, so only gauge-coupling unification (two) survives. The nineteen are free, not derived.

The honest boundary: ``the nineteen are free within the Standard Model'' is a counting fact
(proven); ``no ultraviolet completion derives the values correctly'' is a conjecture --- the
flux theorem worked because the flux has a clean duality group and a first fundamental theorem,
which the Standard Model couplings lack. The surviving, defensible statement is: no unique
\emph{value} is derived --- only \emph{relations}, and the relations that are known (gauge
unification, Yukawa unification) are mostly wrong against data, leaving two.

\section{The universal-structure catalog}

An earlier, stronger guess --- that the sector-spanning content is \emph{one} object (one
K-theory vanishing, one coupling, one observable) --- was tested by three independent literature
passes and \emph{refuted}. What survives is a small catalog:

\begin{enumerate}
\item \textbf{two kinds of universal constraint} --- topological/index vanishing (the tadpole and
the anomaly, the same kind via Atiyah--Singer and anomaly inflow) and positivity/analyticity
(unitarity and causality);
\item \textbf{two universal couplings} --- gravity (the unique coupling to the energy-momentum
tensor) and the dilaton/axion (universal to $\mathrm{Tr}\,F^2$);
\item \textbf{the channel between flux and the Standard Model is the moduli} ---
$1/g^2=\mathrm{Re}\,f(S,T,U)$, the gauge kinetic function, not gravity.
\end{enumerate}

One line: a theory of everything is not one object; it is one \emph{structure} --- constraint,
label, counting --- repeated in every sector, wired together by the moduli.

\section{No selection}

The ``why this vacuum'' question has three levels, and each is answered negatively.

\textbf{Exact.} No invariant selector exists --- this is the theorem of Section~2. \textbf{Statistical.}
The flux measure is flat: measured on $60$ supersymmetric vacua with the physical tadpole bound,
$\arg W_0$ is uniform ($\langle\cos\rangle=0.015$, $\langle\sin\rangle=-0.112$), reproducing the
Denef--Douglas uniformity of $W_0$. No peak, no selection. \textbf{Cosmological.} The pre-uplift
scaling variable $|W_0|^2$ has constant density near zero (measured: the cumulative fraction grows
quadratically), so the cosmological constant is the locally-flat tail of a discretuum --- not a
peak. The uplift that would supply a small positive cosmological constant is \emph{contested}: a
2022--2026 control-problem wave has now hit every uplift mechanism, and the de Sitter swampland is
strongly supported at the ten-dimensional level, though it is not a theorem. The field is trending
toward ``forbidden in the controllable regime,'' with no no-go established.

There is therefore no selection anywhere we can reach: not exact, not statistical, not
cosmological.

\section{The two open gates}

Two items remain open, and they are stated as open, not claimed.

\begin{enumerate}
\item \textbf{The Standard-Model extension (conjectured).} That the nineteen are free in any
ultraviolet completion, beyond the proven within-SM counting, is conjectured. The falsifier is
concrete: any ultraviolet relation among the nineteen beyond GUT's three breaks it.
\item \textbf{The de Sitter uplift (the community's).} Every uplift mechanism now fails the
ten-dimensional control test; the decisive question --- \emph{is the control problem generic?} ---
is string phenomenology's hardest open problem, and this note characterizes it rather than
pretending to close it.
\end{enumerate}

\section{Falsifiers}

The principle is a claim because it can be killed. Four observations would do it.

\begin{enumerate}
\item a moduli-independent, duality-invariant flux quantity that is \emph{not} a function of
$\flux^\top J\hh$ --- breaks the ring $\mathbb{C}[N]$;
\item an exact, derivable, \emph{and numerically-correct} relation among the nineteen beyond
gauge-coupling unification (two) --- breaks the counting;
\item a covariance law that \emph{scatters} at higher $h^{1,2}$ --- breaks the covariance table;
\item a measured non-uniform $dN/d|W_0|^2$ near zero, or a proven de Sitter no-go that empties the
positive-$\Lambda$ discretuum --- breaks the no-selection statement.
\end{enumerate}

\section{Conclusion}

Unification, correctly stated, is not a formula and not one object. It is a complete
characterization: one constraint, a count, a covariance table, and the absence of selection.
That characterization is proven for the flux, measured for the laws and the statistics, and
conjectured --- openly --- for the Standard Model. A theory of everything predicts relations,
never values.

\end{document}
